\section{Conclusion}
In a standardized cohort of 278 human pelves evaluated under finite-element loading, sacroiliac joint micromotion revealed four fundamental biomechanical principles:
\begin{enumerate}
\item \textbf{Load-Dependence and Stance Asymmetry:} Erect bipedal standing reinforces form and force closure through physiological nutation ($\theta_{\text{nut}} = 1.062^\circ$), whereas single-leg stance imposes the highest kinematic demand among the evaluated regimes, driving a marked increase in bilateral translation asymmetry ($2.013$~mm vs.\ $0.215$~mm) and frontal tilt.
\item \textbf{Directional Kinematic Contrasts:} Parturition-motivated forces suppress sagittal nutation and vertical shear, shifting joint displacement toward transverse outward distraction ($\Delta\text{ML} \ge +0.09$~mm) and sagittal gliding, illustrating passive compliance pathways that facilitate cavity expansion while maintaining articular congruency.
\item \textbf{Sex-Related Differences and Architectural Associations:} Female pelves exhibit significantly greater translational compliance across labor proxies ($+0.19$ to $+0.23$~mm in fully adjusted models, $p \le 0.020$). This difference attenuates upon adjusting for pelvic volume and dimensions, and correlates positively with subpubic arch angle ($r = 0.52, p < 10^{-19}$) due to the lower hoop stiffness of broader anterior arch morphology.
\item \textbf{Dominance of Skeletal Geometry over Bone Mineralization Heterogeneity:} Macro-architectural pelvic anatomy accounts for 98.7\% to 102.1\% (median 100.5\%) of between-subject kinematic variance ($R^2 \ge 0.994$), whereas bone mineral density heterogeneity explains less than 0.32\% (median 0.033\%), establishing that 3D skeletal geometry, rather than bone mineralization, governs joint-level kinematic boundaries.
\end{enumerate}
Together, these findings demonstrate that sacroiliac joint kinematics are governed by the interplay of load configuration and skeletal architecture, clarifying the mechanical compromise between bipedal stability and pelvic compliance, and providing quantitative anatomical benchmarks for orthopedic reconstruction and pelvic biomechanics.
